\section{Technical Details}

The model was created using pytorch in the Google Collab Lab environment. 
The dataset was constructed in memory and then a logical AND was applied to create labels.
The ReLU~\ref{fig:relu-eq} was used for the activation function and MSE~\ref{fig:mse-loss-eq} was used for the loss function.
Stochastic Gradient Decent was used for optimization.
Training was performed with a learning rate of 0.03, 100 epochs, and a loss threshold of 0.10.


\begin{figure}[ht]
  \centering
  \begin{equation}
    \mathrm{ReLU}(x) = \max(0, x)
  \end{equation}
  \caption{Definition of the ReLU activation function}
  \label{fig:relu-eq}
\end{figure}

\begin{figure}[ht]
  \centering
  \begin{equation}
    \mathrm{MSE}(y, \hat{y}) =
    \frac{1}{n} \sum_{i=1}^{n} \bigl(y_i - \hat{y}_i\bigr)^2
  \end{equation}
  \caption{Mean Squared Error (MSE) loss function}
  \label{fig:mse-loss-eq}
\end{figure}

Training was performed using forward propagation. The below formulas demonstrate the forward propagation matrix operations.
The weights are multiplied with the input features and a bias is added. A activation function is then applied to that operation to introduce non-linearity.
Each subsequent layer follows this process using the previous activations.
The output layer creates a prediction. The activation here can be different than previous layers. For this model ReLU was utilized.


\begin{align}
z^{(1)} &= W^{(1)} x + b^{(1)}, \\
a^{(1)} &= \sigma\!\bigl(z^{(1)}\bigr), \\
z^{(2)} &= W^{(2)} a^{(1)} + b^{(2)}, \\
\hat{y} &= f\!\bigl(z^{(2)}\bigr),
\end{align}

The model weights and bias' were initialized with a random distribution. 
Training was performed in a loop and terminated on a epoch or a loss threshold.
If the loss dropped below a stated threshold, training ceases to save time.
The epoch termination is for when the model converges.
Training only on a loss threshold can create a infinite loop if the convergence is higher than the threshold.

\begin{figure}[htpb]
    \centering
    \includegraphics[width=0.8\linewidth]{./figures/x_tensor.png}
    \caption{Screen shot of initialization of weights and bias', dataset and label creation, and model organization.}
    \label{fig:init}
\end{figure}

\begin{figure}[htpb]
\centering
\includegraphics[width=0.8\linewidth]{./figures/training.png}
    \caption{Screen shot of training loop and print statements.}
\label{fig:training}
\end{figure}

